%% kapa: exam=1 round=32 year=2021 period=1 form=A subject=economics subject_ko=경제학원론 seq=15 number=55
%% source: Q-Net 감정평가사 기출문제 (cst00309) artlSeq=5208397 file=extracted/1008/5208397/A형/2021년도 제32회 감정평가사 1차 1교시 A형.pdf
%% conversion: kordoc 4.24.51 -> tools/gen1.py (auto); review=proofread-sonnet; flags=none
\begin{question}{55}
순수 공공재에 관한 설명으로 옳지 않은 것은?
\begin{choices}
  \item 소비자가 많을수록 개별 소비자가 이용하는 편익은 감소한다.
  \item 시장수요는 개별 소비자 수요의 수직합으로 도출된다.
  \item 개별 소비자의 한계편익 합계와 공급에 따른 한계비용이 일치하는 수준에서 사회적 최적량이 결정된다.
  \item 시장에서 공급량이 결정되면 사회적 최적량에 비해 과소 공급된다.
  \item 공급량이 사회적 최적 수준에서 결정되려면 사회 전체의 정확한 선호를 파악해야 한다.
\end{choices}
\end{question}
\answer{1}
